\section{Exact level counts and the native dynamical zeta function}
\label{sec:counts}

Fix integers $a,k$. Let $c_\ell(a,k)$ be the number of
oriented cycles of least period $\ell$ on the ordinary integer
level set
\[
 X_{a,k}(\Z)=\{(x,y,z)\in\Z^3:K_a(x,y,z)=k\}.
\]
All lengths outside \eqref{eq:periodset} have count zero.
The following explicit tests remove every remaining orbit search.
Here a positive square means $h^2$ for an integer $h>0$;
zero is excluded whenever positivity is specified.

\begin{proposition}[Finite arithmetic tests]
\label{prop:counts}
The counts $c_\ell(a,k)$ are given by the following rules.
\begin{enumerate}
\item $c_1$ is the number of distinct integers
$r=1\pm\sqrt{1-a}$ satisfying $k=r^2(2r-3)$.
If $1-a$ is not a nonnegative integer square, the count is
zero; if the square root is zero, count the common root once.
\item Let $P$ run over the distinct integer roots of
\begin{equation}
 P^2+(a-3)P-k=0.
 \label{eq:period2quadratic}
\end{equation}
Put $S=P+a$. Each root contributes one to $c_2$ exactly
when $S^2-4P=h^2>0$ with $h\equiv S\pmod2$.
Otherwise it contributes zero.
\item $c_3=1$ exactly when $a$ is even,
$t=(a+4)/2\ne-2$, and $k=8-(t-4)^2$; otherwise $c_3=0$.
\item $c_4=1$ exactly when
\begin{equation}
 (a+1)(a-7)+4k=h^2>0,\qquad h\equiv a+1\pmod2;
 \label{eq:period4square}
\end{equation}
otherwise $c_4=0$.
\item At $a=-1$, the $F_5$ contribution is one if $k=0$,
two if $k>0$ and $4k+1$ is an integer square, and zero
otherwise. Add one precisely for $(a,k)=(-13,-64)$.
There is no other contribution to $c_5$.
\item $c_6=1$ exactly when $a=0$ and $k$ is a positive
integer square.
\item $c_8=1$ exactly when $a=-1$ and $k-1$ is a
positive integer square.
\item $c_9=1$ exactly when $(a,k)=(-2,-1)$.
\item $c_{12}=1$ exactly when $a=0$ and $4k-7$ is a
positive integer square.
\end{enumerate}
In the last four rules the count is zero when the stated
condition fails.
\end{proposition}

\begin{proof}
The fixed-point relation is $(r-1)^2=1-a$, giving the
first test. For an alternating pair put $S=u+v$ and
$P=uv$. The parameter and level equations become
\[
 a=S-P,\qquad k=P(S-3)=P(P+a-3).
\]
This proves \eqref{eq:period2quadratic}. A distinct integral
unordered pair exists exactly when its discriminant is a
positive square with the required parity; then
\[
 u=(S-h)/2<v=(S+h)/2.
\]
Distinct values of $P$ give distinct pairs, since $S=P+a$.
This formulation remains valid at $a=1$ and does not divide
by $u-1$ or lose a factor-zero case.

The $F_3$ relation uniquely fixes $t$ when $a$ is even.
For $F_4$, solving
$k=t^2-(a+1)t+2a+2$ gives discriminant
$(a+1)(a-7)+4k$. Its smaller root is
$(a+1-h)/2$; it is integral and satisfies $2t<a+1$
exactly under \eqref{eq:period4square}. Thus there is
one cycle, not two, when that test succeeds.

For $F_5$, $k=t(t+1)\ge0$ for integral $t$.
At $k=0$ the two roots $0,-1$ give the same cyclic
word and the specified range retains only $0$. At $k>0$
the square test $4k+1=(2t+1)^2$ gives two distinct
integer parameters, neither excluded; they give the two
distinct reversed oriented cycles and must both be counted.
The additional exceptional contribution is read at its
separate parameter and level.

The formulas $k=t^2$ with $t\ge1$ and
$k-1=(t+1)^2$ with $t\ge0$ give the sixth and
seventh tests. The eighth is the unique $E_9$ row.
Finally $k=m^2-m+2$ is equivalent to
$4k-7=(2m-1)^2$. A positive square congruent to $1$
modulo $4$ has odd positive root $h$, yielding the unique
$m=(1+h)/2\ge1$. The disjointness already proved means
that no further identifications enter these tests.
\end{proof}

\begin{corollary}[Ordinary fixed points and zeta]
\label{cor:zeta}
For all integers $a,k$ and $n\ge1$,
\begin{equation}
 \#\Fix\bigl(T_a^n\vert X_{a,k}(\Z)\bigr)
 =\sum_{\substack{\ell\in\cL\\\ell\mid n}}
   \ell c_\ell(a,k).
 \label{eq:fixedcounts}
\end{equation}
The Artin--Mazur dynamical zeta series in the ordinary time
variable $\tau$ is the finite rational product
\begin{equation}
 \zeta_{a,k}(\tau)
 :=\exp\!\left(\sum_{n\ge1}
 \frac{\#\Fix(T_a^n\vert X_{a,k}(\Z))}{n}\tau^n\right)
 =\prod_{\ell\in\cL}(1-\tau^\ell)^{-c_\ell(a,k)}.
 \label{eq:zeta}
\end{equation}
\end{corollary}
\begin{proof}
Each arithmetic test admits at most finitely many parameters,
so $X_{a,k}(\Z)$ has finitely many periodic points.
An oriented cycle of length $\ell$ contributes its
$\ell$ distinct points to $\Fix(T_a^n)$ exactly when
$\ell\mid n$, proving \eqref{eq:fixedcounts}.
For each such cycle its contribution to the logarithm of
the zeta series is
$\sum_{j\ge1}\tau^{j\ell}/j=-\log(1-\tau^\ell)$.
Summing over the finitely many cycles proves
\eqref{eq:zeta} as a formal power-series identity.
\end{proof}

There is no analogous unrestricted all-level finite-count
assertion here. For every fixed $a$, the infinitely many
integers $t$ with $2t<a+1$ give distinct four-cycles.
In particular the all-level fixed set of $T_a^4$ is
infinite. Fixing the invariant level is indispensable for
the finite counts in \eqref{eq:fixedcounts}.
